% Options for packages loaded elsewhere
\PassOptionsToPackage{unicode}{hyperref}
\PassOptionsToPackage{hyphens}{url}
\documentclass[
  ignorenonframetext,
]{beamer}
\newif\ifbibliography
\usepackage{pgfpages}
\setbeamertemplate{caption}[numbered]
\setbeamertemplate{caption label separator}{: }
\setbeamercolor{caption name}{fg=normal text.fg}
\beamertemplatenavigationsymbolsempty
% remove section numbering
\setbeamertemplate{part page}{
  \centering
  \begin{beamercolorbox}[sep=16pt,center]{part title}
    \usebeamerfont{part title}\insertpart\par
  \end{beamercolorbox}
}
\setbeamertemplate{section page}{
  \centering
  \begin{beamercolorbox}[sep=12pt,center]{section title}
    \usebeamerfont{section title}\insertsection\par
  \end{beamercolorbox}
}
\setbeamertemplate{subsection page}{
  \centering
  \begin{beamercolorbox}[sep=8pt,center]{subsection title}
    \usebeamerfont{subsection title}\insertsubsection\par
  \end{beamercolorbox}
}
% Prevent slide breaks in the middle of a paragraph
\widowpenalties 1 10000
\raggedbottom
\AtBeginPart{
  \frame{\partpage}
}
\AtBeginSection{
  \ifbibliography
  \else
    \frame{\sectionpage}
  \fi
}
\AtBeginSubsection{
  \frame{\subsectionpage}
}
\usepackage{iftex}
\ifPDFTeX
  \usepackage[T1]{fontenc}
  \usepackage[utf8]{inputenc}
  \usepackage{textcomp} % provide euro and other symbols
\else % if luatex or xetex
  \usepackage{unicode-math} % this also loads fontspec
  \defaultfontfeatures{Scale=MatchLowercase}
  \defaultfontfeatures[\rmfamily]{Ligatures=TeX,Scale=1}
\fi
\usepackage{lmodern}
\ifPDFTeX\else
  % xetex/luatex font selection
\fi
% Use upquote if available, for straight quotes in verbatim environments
\IfFileExists{upquote.sty}{\usepackage{upquote}}{}
\IfFileExists{microtype.sty}{% use microtype if available
  \usepackage[]{microtype}
  \UseMicrotypeSet[protrusion]{basicmath} % disable protrusion for tt fonts
}{}
\makeatletter
\@ifundefined{KOMAClassName}{% if non-KOMA class
  \IfFileExists{parskip.sty}{%
    \usepackage{parskip}
  }{% else
    \setlength{\parindent}{0pt}
    \setlength{\parskip}{6pt plus 2pt minus 1pt}}
}{% if KOMA class
  \KOMAoptions{parskip=half}}
\makeatother
\setlength{\emergencystretch}{3em} % prevent overfull lines
\providecommand{\tightlist}{%
  \setlength{\itemsep}{0pt}\setlength{\parskip}{0pt}}
\usepackage{bookmark}
\IfFileExists{xurl.sty}{\usepackage{xurl}}{} % add URL line breaks if available
\urlstyle{same}
\hypersetup{
  hidelinks,
  pdfcreator={LaTeX via pandoc}}

\author{\texorpdfstring{}{}}
\date{}

\begin{document}

\begin{frame}{Consolidating Fiat Abstraction: Federal Reserve to the
Nixon Shock}
\protect\phantomsection\label{consolidating-fiat-abstraction-federal-reserve-to-the-nixon-shock}
The arc from Newton's 1717 ratio through Hamilton's 1792 ratio traces a
single structural trajectory: the global establishment of monometallic
quantification over bimetallic generation. Yet the 19th-century gold
standard was merely a transitional phase toward a more radical
abstraction. The defining monetary events of the 20th century---the
establishment of the Federal Reserve (1913), the confiscation of private
gold (1933), the demonetization of silver coinage (1965), and the
collapse of the Bretton Woods system (1971)---constitute the realization
of Blake's \emph{Ulro}: the descent into a fiat void where money is
severed entirely from physical substance.

\begin{block}{The 1913 Federal Reserve Act: Centralizing the
Architecture of Issuance}
\protect\phantomsection\label{the-1913-federal-reserve-act-centralizing-the-architecture-of-issuance}
On 23 December 1913, Congress passed the Federal Reserve Act, altering
the architecture of American monetary sovereignty by transferring the
power of currency issuance to a centralized, quasi-public banking system
comprising twelve regional Reserve Banks coordinated by the Federal
Reserve Board in Washington, D.C. As William Silber documents in
\emph{When Washington Shut Down Wall Street}, the establishment of the
Federal Reserve was catalyzed not merely by the Panic of 1907 but by a
structural crisis: the absence of a ``lender of last resort'' capable of
providing elastic currency during financial panics---a function the Bank
of England had performed since 1694 {[}@silber2007when{]}. The Federal
Reserve thus represented the belated American adoption of the
institutional architecture that Blake had identified as the engine of
abstract credit creation.

In Blake's mythological system, the Federal Reserve Act represents the
institutional triumph of the ``Councellor'' who invents, through a kind
of fiat, ``allegoric riches'' (\emph{Jerusalem} 27; see
\hyperlink{sec-allegoric-riches}{§ The fixing of labour and allegoric
riches}) {[}@erdman1988complete{]}. By establishing an elastic currency
supply managed through open market operations and fractional reserve
banking, the Federal Reserve decoupled the dollar from any fixed
metallic constraint, allowing the money supply to expand and contract at
the discretion of an appointed technocratic elite. Liaquat Ahamed's
\emph{Lords of Finance} demonstrates how the central bankers of the
interwar period---Benjamin Strong of the New York Fed, Montagu Norman of
the Bank of England, Émile Moreau of the Banque de France, and Hjalmar
Schacht of the Reichsbank---wielded this discretionary power with severe
economic consequences, their collective mismanagement of the gold
exchange standard directly precipitating the Great Depression
{[}@ahamed2009lords{]}.

A reading of Blake's Four Zoas as a ``proto-cognitive architecture''
illuminates the Federal Reserve's structural error: the consolidation of
monetary authority in a single faculty---Urizen's reason---detached from
the embodied, sensory, and imaginative dimensions that the other Zoas
represent {[}@friedman2026blake{]}. The Fed's founding premise---that
rational technocratic management can optimize the money supply---is
precisely the Urizenic fantasy that Blake diagnosed. It operationalizes
``Newton's sleep,'' which functions formally as the pathology of rigid
priors crushing sensory evidence: the belief that a single geometric
framework, administered by a priestly elite, can govern the infinitely
complex organism of economic life {[}@friedman2026doors{]}. Where the
earlier bimetallic tension decomposed value into two competing metrics,
the fiat regime represents an even deeper loss of structural integrity.
Synergetic wholeness is destroyed when value becomes an arbitrarily
manipulable abstract symbol, severing any remaining anchoring to
physical reality and allowing limitless numerical expansion without
commensurate productive generation.
\end{block}

\begin{block}{Executive Order 6102: The Confiscation of the Solar
Principle}
\protect\phantomsection\label{executive-order-6102-the-confiscation-of-the-solar-principle}
If the Federal Reserve Act abstracted the issuance of money, the onset
of the Great Depression catalyzed the confiscation of the metallic
anchor itself. Under Executive Order 6102 (5 April 1933) and the
subsequent Gold Reserve Act of 30 January 1934, President Franklin D.
Roosevelt mandated the surrender of all privately held gold coins,
bullion, and gold certificates to the Federal Reserve
{[}@roosevelt1933executive{]}. Private gold ownership was
criminalized---punishable by up to ten years' imprisonment and a
\$10,000 fine---and the metal was collected and revalued from \$20.67 to
\$35.00 per ounce. This 69\% devaluation effectively expropriated the
savings of every American who had held wealth in the constitutionally
mandated medium of exchange, while simultaneously concentrating the
nation's gold in federal vaults at Fort Knox (constructed 1936) and West
Point.

The ``mental chains'' described in \emph{America a Prophecy} had thus
materialized into statutory edicts: the citizen could no longer hold the
physical substance of value. By capturing the gold, the federal
government removed a physical boundary against fiat expansion. The
contemporaneous Silver Purchase Act of 1934, championed by the Western
silver bloc in Congress, attempted to inflate the currency by
nationalizing vast silver reserves and mandating Treasury purchases at
artificially elevated prices. This dual confiscation demonstrates that
the state recognized the necessity of capturing \emph{both} poles of the
bimetallic spectrum to enforce systemic compliance---an institutional
consolidation that Blake's prophetic framework renders legible as the
final enclosure of metallic reality within the Urizenic ledger.
\end{block}

\begin{block}{The 1965 Coinage Act and the Purging of Circulating
Silver}
\protect\phantomsection\label{the-1965-coinage-act-and-the-purging-of-circulating-silver}
This trajectory reached its domestic culmination with the Coinage Act of
23 July 1965, signed by President Lyndon B. Johnson. Driven by acute
shortages resulting from the growing divergence between silver's
commodity value and its face value in coinage---a textbook Gresham's Law
phenomenon in which the metal's intrinsic worth exceeded its nominal
denomination---the Act removed 90\% silver content from circulating
dimes and quarters, replacing them with copper-nickel clad tokens
{[}@johnson1965coinage{]}. Half dollars were reduced from 90\% to 40\%
silver, and would lose their remaining silver content entirely in 1971.
Johnson acknowledged the magnitude of this transformation in his signing
remarks, assuring Congress that the Treasury's silver reserves would
``keep the price of silver in line with its value in our present silver
coin.''

The purging of silver from everyday exchange completed a process that
had begun with the Crime of 1873. What circulated was no longer physical
wealth but symbolic tokens, dependent entirely on the credibility of the
issuing authority. The United States had achieved domestically what
Blake diagnosed as the terminal condition of Ulro: the total replacement
of intrinsic material substance with abstract, administered fiat.
\end{block}

\begin{block}{The Nixon Shock and the Global Triumph of Pure
Abstraction}
\protect\phantomsection\label{the-nixon-shock-and-the-global-triumph-of-pure-abstraction}
The ontological transformation of money was achieved globally on 15
August 1971, when President Richard Nixon unilaterally terminated the
direct international convertibility of the US dollar to gold.
Confronting the depletion of US gold reserves from a postwar peak of
approximately 20,000 tonnes to roughly 8,133 tonnes, alongside the
deficit financing of the Vietnam War and Great Society programs, Nixon
suspended the gold window that had anchored the post-war Bretton Woods
system since 1944 {[}@steil2013battle; @eichengreen2019globalizing{]}.

Benn Steil's \emph{The Battle of Bretton Woods} documents how Keynes and
White had designed the original system as a managed gold exchange
standard precisely to avoid the deflationary rigidity of the classical
gold standard while preserving some metallic discipline against
inflationary excess {[}@steil2013battle{]}. Nixon's unilateral
abrogation of this architecture---announced on a Sunday evening without
consulting allied governments---severed the final tether between the
global monetary system and physical reality. The world was plunged into
a purely fiat regime---a system where all global currencies float
against one another, sustained entirely by the credibility, military
power, and institutional inertia of sovereign issuers. As Barry
Eichengreen demonstrates in \emph{Globalizing Capital}, no historical
precedent existed for a global fiat system of this scale; every previous
experiment with unbacked paper money, from the Chinese \emph{jiaozi} of
the Song Dynasty to the French \emph{assignats} of the 1790s, had ended
in inflationary collapse {[}@eichengreen2019globalizing{]}.

Official US Treasury--owned monetary gold is often reported near 8,133
tonnes (post--Bretton Woods levels; exact totals move with coining and
transfers). Site figures for Fort Knox, West Point, Denver, and related
facilities are components of that sovereign stock, not separate pools to
be summed into a larger national total. Metal vaulted at the Federal
Reserve Bank of New York is predominantly foreign and other custodial
holdings---not US official reserves---so it must not be added to
Treasury aggregates. That accounting split between sovereign stock and
Manhattan custody mirrors the ontological crisis Blake diagnosed: the
ledger's line items and the underlying metal are easy to conflate.
Treasury gold remains officially valued on the books at the statutory
price of \$42.22 per ounce established in 1973, vastly understating
contemporary market reality. Policy discussion since 2025 has examined
historical precedents for sovereign gold revaluations, estimating that
marking US reserves to market would yield proceeds equivalent to roughly
3\% of GDP---a figure that has prompted Congressional interest in
revaluation as a mechanism for sovereign balance sheet management. The
Urizenic ledger maintains its geometric fiction---\$42.22---while the
physical metal trades at a far higher market price.
\end{block}
\end{frame}

\end{document}
